% !PRINTED 250
\weberpage{250}
Daraus ergibt sich wie oben der Schluß:

1. \gesperrt{Für jedes beliebige $n$ sind die Größen}
\[
P^{24}_{c,\eth,a}
\]
\gesperrt{die Wurzeln einer Transformationsgleichung.}

Ist $n$ durch 3 nicht teilbar, so kann man die $c, c_2, c_1$ durch
3 teilbar voraussetzen. Dann wird
\[
\lambda \equiv n, \quad \alpha \equiv 0, \quad \delta \equiv 0 \pmod{3}.
\]
Es ist also, wie aus §\,38, (15) hervorgeht, $r$ eine \gesperrt{achte
Einheitswurzel}, beachtet man daher noch
\[
\gamma_2(\omega+1) = e^{-\frac{2\pi i}{3}}\gamma_2(\omega), \quad
\gamma_2\left(-\frac{1}{\omega}\right) = \gamma_2(\omega), \quad \text{[§\,54, (14)]}
\]
so haben wir den Satz:

2. \gesperrt{Ist $n$ nicht durch 3 teilbar und $c$ durch 3 teilbar,
so sind die Größen}
\[
P^{8}_{c,\eth,a}\,\gamma_2(\omega)^{n-1}
\]
\gesperrt{Wurzeln einer Transformationsgleichung.}

Um die Anwendung auf den einfachsten Fall $n=2$ zu machen, setzen wir

\weq{3}{x = 16\,\frac{\eta(2\omega)^8}{\eta(\omega)^8}, \quad
x_0 = \frac{\eta\left(\dfrac{\omega}{2}\right)^8}{\eta(\omega)^8}, \quad
x_1 = \frac{\eta\left(\dfrac{\omega+3}{2}\right)^8}{\eta(\omega)^8},}

und dann sind $x, x_0, x_1$ nichts anderes als unsere Funktionen
\[
f_2(\omega)^8, \quad f_1(\omega)^8, \quad -f(\omega)^8. \qquad \text{[§\,34, (9)]}
\]
Diese sind, wie wir schon früher gesehen haben [§\,54, (8)], die
Wurzeln der kubischen Gleichung
\weq{4}{x^3 - x\,\gamma_2(\omega) + 16 = 0.}
Der Umstand, daß $f_2^8, f_1^8, -f^8$ selbst Wurzeln einer Transformationsgleichung für den Transformationsgrad 2 sind, erklärt die
Erscheinung, daß bei Adjunktion dieser Größen, also auch bei
Adjunktion von $x^2$, die zu einem geraden $n$ gehörigen Transformationsgleichungen reducibel werden.

Wenn $n$ eine ungerade Zahl ist, so nehme man $c$ \gesperrt{durch 4 teilbar}
an, dann ist [§\,69, (9), (11), (12)]
\[
\lambda \equiv n \pmod{4},
\]
\[
\alpha \equiv 0, \quad \delta \equiv 0, \quad \beta \equiv a\,\eth_2, \quad \gamma \equiv -\eth\,a_2 \pmod{4},
\]
also ergibt sich
\[
r^6 = (-1)^{\frac{\beta-1}{2}} = (-1)^{\frac{n-1}{2}}(-1)^{\frac{a-a_2}{2}}
\]
%CONT>>
